Python 的动态行为，以及默认情况下缺乏静态类型检查（static type checking），既是一种恩赐，也是一种诅咒（详见 Item 3「绝不要期望 Python 在编译期检测错误（Never Expect Python to Detect Errors at Compile Time）」）。然而，数量庞大的 Python 程序员表示，它值得使用，因为由此带来的简洁与单纯带来了生产率。但大多数使用 Python 的人，都至少有一个关于程序在运行时撞上愚蠢错误的恐怖故事。

我听过的最糟糕的例子之一，涉及一个 \texttt{SyntaxError} 异常因一次动态导入（dynamic import）的副作用而在生产环境中被抛出（参见 Item 98「用动态导入惰性加载模块以减少启动时间（Lazy-Load Modules with Dynamic Imports to Reduce Startup Time）」的例子），导致服务器进程崩溃。

我所认识的这位被这桩意外事件击中的程序员，自此将使用 Python 的可能性彻底排除。但我不得不怀疑：在程序部署到生产环境之前，这段代码为何没有经过更充分的测试？编译期静态类型安全并非一切。无论代码用什么语言编写，你都应当始终测试它。不过我承认，在 Python 中，编写测试来验证正确性，可能比其他语言更为重要。

所幸，造成这些风险的同样一些动态特性，也让为你的代码编写测试、调试故障程序变得极其容易。你可以利用 Python 的动态本性，以及易于覆盖的行为，来实现测试并确保你的程序按预期工作。

\section{Item 108: 在 TestCase 子类中验证相关行为}

在 Python 中编写测试的规范做法，是使用内置（built-in）模块 \texttt{unittest}。例如，假设我定义了如下这个工具函数，我希望验证它在各种输入下都能正确工作：

\begin{codebox}[python]
# utils.py
def to_str(data):
    if isinstance(data, str):
        return data
    elif isinstance(data, bytes):
        return data.decode("utf-8")
    else:
         raise TypeError(
    f"Must supply str or bytes, found: {data}")
\end{codebox}

为了定义测试，我会创建第二个文件，命名为 \texttt{test\_utils.py} 或 \texttt{utils\_test.py}——具体命名方式是风格选择——其中包含我期望的每种行为的测试：

\begin{codebox}[python]
# utils_test.py
from unittest import TestCase, main
from utils import to_str
class UtilsTestCase(TestCase):
    def test_to_str_bytes(self):
        self.assertEqual("hello", to_str(b"hello"))
    def test_to_str_str(self):
        self.assertEqual("hello", to_str("hello"))
    def test_failing(self):
        self.assertEqual("incorrect", to_str("hello"))
if __name__ == "__main__":
    main()
\end{codebox}

然后我使用 Python 命令行运行这个测试文件。在本例中，三个测试方法中两个通过、一个失败，并打印出有助于定位问题的错误消息：

\begin{codebox}[python]
$ python3 utils_test.py
F..
===============================================================
FAIL: test_failing (__main__.UtilsTestCase)
---------------------------------------------------------------
Traceback (most recent call last):
  File "utils_test.py", line 15, in test_failing
    self.assertEqual('incorrect', to_str('hello'))
AssertionError: 'incorrect' != 'hello'
- incorrect
+ hello
---------------------------------------------------------------
Ran 3 tests in 0.002s
FAILED (failures=1)
\end{codebox}

测试被组织进 \texttt{TestCase} 子类中。每个测试用例都是一个以 \texttt{test} 开头的方法。如果某个测试方法在运行时没有抛出任何异常（包括来自 \texttt{assert} 语句的 \texttt{AssertionError}；参见 Item 81：\texttt{assert} 内部假设，\texttt{raise} 未达预期），那么该测试就被认为已通过。如果某个测试失败，\texttt{TestCase} 子类会继续运行其他测试方法，这样你就能掌握所有测试的整体情况，而不是在第一个问题出现时就停下来。

如果你想快速迭代来修复或改进某个特定测试，可以在命令行上指定该测试方法在测试模块中的路径，从而只运行那一个测试方法：

\begin{codebox}[python]
$ python3 utils_test.py UtilsTestCase.test_to_str_bytes
.
---------------------------------------------------------------
Ran 1 test in 0.000s
OK
\end{codebox}

你也可以直接在测试方法内部、在特定的断点处调用调试器，以更深入地排查失败的原因（具体做法参见 Item 114：考虑使用 \texttt{pdb} 进行交互式调试）。

\texttt{TestCase} 类提供了用于在测试中做断言的辅助方法，例如用于验证相等的 \texttt{assertEqual}、用于验证布尔表达式的 \texttt{assertTrue}、在关注精度时使用的 \texttt{assertAlmostEqual}（参见 Item 113：使用 \texttt{assertAlmostEqual} 控制浮点测试中的精度），以及其他许多方法（完整列表见 https://docs.python.org/3/library/unittest.html）。它们比内置的 \texttt{assert} 语句更好，因为它们会打印出所有输入与输出，帮助你准确理解测试失败的原因。例如，下面我分别用辅助断言方法和不用辅助断言方法写了同一个测试用例：

\begin{codebox}[python]
# assert_test.py
from unittest import TestCase, main
from utils import to_str
class AssertTestCase(TestCase):
    def test_assert_helper(self):
        expected = 12
        found = 2 * 5
        self.assertEqual(expected, found)
    def test_assert_statement(self):
        expected = 12
        found = 2 * 5
        assert expected == found
if __name__ == "__main__":
    main()
\end{codebox}

这两条失败消息，哪一条对你更有帮助？注意第二条消息没有显示 \texttt{expected} 或 \texttt{found} 的值：

\begin{codebox}[python]
$ python3 assert_test.py
FF
===============================================================
FAIL: test_assert_helper (__main__.AssertTestCase)
---------------------------------------------------------------
Traceback (most recent call last):
  File "assert_test.py", line 16, in test_assert_helper
    self.assertEqual(expected, found)
AssertionError: 12 != 10
===============================================================
FAIL: test_assert_statement (__main__.AssertTestCase)
---------------------------------------------------------------
Traceback (most recent call last):
  File "assert_test.py", line 11, in test_assert_statement
    assert expected == found
AssertionError
---------------------------------------------------------------
Ran 2 tests in 0.001s
FAILED (failures=2)
\end{codebox}

还有一个 \texttt{assertRaises} 辅助方法，用于验证异常；它可以在 \texttt{with} 语句中作为上下文管理器使用（其工作方式参见 Item 82：考虑用 \texttt{contextlib} 与 \texttt{with} 语句实现可复用的 \texttt{try}/\texttt{finally} 行为）。它看起来类似于 \texttt{try}/\texttt{except} 语句，并且能清晰地表明应该在何处抛出该异常：

\begin{codebox}[python]
# utils_error_test.py
from unittest import TestCase, main
from utils import to_str
class UtilsErrorTestCase(TestCase):
    def test_to_str_bad(self):
        with self.assertRaises(TypeError):
            to_str(object())
    def test_to_str_bad_encoding(self):
        with self.assertRaises(UnicodeDecodeError):
            to_str(b"\xfa\xfa")
if __name__ == "__main__":
    main()
\end{codebox}

你还可以在 \texttt{TestCase} 子类中定义带有复杂逻辑的自有辅助方法，使测试更易读。只要确保方法名不以 \texttt{test} 开头，否则它们会被当作测试用例来运行。除了调用 \texttt{TestCase} 的断言方法，这些自定义测试辅助方法经常使用 \texttt{fail} 方法来说明是哪个假设或不变量（invariant）没有被满足。例如，这里我定义了一个自定义测试辅助方法，用于验证生成器（generator）的行为：

\begin{codebox}[python]
# helper_test.py
from unittest import TestCase, main
def sum_squares(values):
    cumulative = 0
    for value in values:
        cumulative += value**2
        yield cumulative
class HelperTestCase(TestCase):
    def verify_complex_case(self, values, expected):
        expect_it = iter(expected)
        found_it = iter(sum_squares(values))
        test_it = zip(expect_it, found_it, strict=True)
        for i, (expect, found) in enumerate(test_it):
            if found != expect:
                self.fail( f”Index {i} is wrong: “
f”{found} != {expect}”)
    def test_too_short(self):
        values = [1.1, 2.2]
        expected = [1.1**2]
        self.verify_complex_case(values, expected)
    def test_too_long(self):
        values = [1.1, 2.2]
        expected = [
            1.1**2,
            1.1**2 + 2.2**2,
            0,  # Value doesn't matter
        ]
        self.verify_complex_case(values, expected)
    def test_wrong_results(self):
        values = [1.1, 2.2, 3.3]
        expected = [
            1.1**2,
            1.1**2 + 2.2**2,
            1.1**2 + 2.2**2 + 3.3**2 + 4.4**2,
        ]
        self.verify_complex_case(values, expected)
if __name__ == "__main__":
    main()
\end{codebox}

这个辅助方法让测试用例变得简洁、易读，输出的错误消息也容易理解：

\begin{codebox}[python]
$ python3 helper_test.py
EEF
=============================================================
ERROR: test_too_long (__main__.HelperTestCase.test_too_long)
-------------------------------------------------------------
Traceback (most recent call last):
  File "helper_test.py", line 36, in test_too_long
    self.verify_complex_case(values, expected)
    ~~~~~~~~~~~~~~~~~~~~~~~~^^^^^^^^^^^^^^^^^^
  File "helper_test.py", line 20, in verify_complex_case
    for i, (expect, found) in enumerate(test_it):
                              ~~~~~~~~~^^^^^^^^^
ValueError: zip() argument 2 is shorter than argument 1
=============================================================
ERROR: test_too_short (__main__.HelperTestCase.test_too_short)
-------------------------------------------------------------
Traceback (most recent call last):
  File "helper_test.py", line 27, in test_too_short
    self.verify_complex_case(values, expected)
    ~~~~~~~~~~~~~~~~~~~~~~~~^^^^^^^^^^^^^^^^^^
  File "helper_test.py", line 20, in verify_complex_case
    for i, (expect, found) in enumerate(test_it):
                              ~~~~~~~~~^^^^^^^^^
ValueError: zip() argument 2 is longer than argument 1
=============================================================
FAIL: test_wrong_results
'(__main__.HelperTestCase.test_wrong_results)
-------------------------------------------------------------
Traceback (most recent call last):
  File "helper_test.py", line 45, in test_wrong_results
    self.verify_complex_case(values, expected)
    ~~~~~~~~~~~~~~~~~~~~~~~~^^^^^^^^^^^^^^^^^^
  File "helper_test.py", line 22, in verify_complex_case
    self.fail(f"Index {i} is wrong: {found} != {expect}")
    ~~~~~~~~~^^^^^^^^^^^^^^^^^^^^^^^^^^^^^^^^^^^^^^^^^^^^
AssertionError: Index 2 is wrong: 16.939999999999998 != 36.3
-------------------------------------------------------------
Ran 3 tests in 0.001s
FAILED (failures=1, errors=2)
\end{codebox}

我通常会为每一组相关的测试定义一个 \texttt{TestCase} 子类。有时，我会为每个拥有很多边界情况的函数定义一个 \texttt{TestCase} 子类；另一些时候，一个 \texttt{TestCase} 子类会覆盖单个模块里的所有函数。我经常为每个基础类及其所有方法定义一个 \texttt{TestCase} 子类来做测试（更多指引参见 Item 109：优先选择集成测试而非单元测试）。

\texttt{TestCase} 类还提供了一个 \texttt{subTest} 辅助方法，让你得以在一个测试方法内定义多个测试，从而避免样板代码。它在编写数据驱动（data-driven）测试时尤其有用，并且即使其中某个用例失败，测试方法仍能继续测试其他用例（这与 \texttt{TestCase} 对待其包含的各个测试方法的行为类似；另一种做法参见 Item 110：用 \texttt{setUp}、\texttt{tearDown}、\texttt{setUpModule} 和 \texttt{tearDownModule} 将测试彼此隔离）。为了说明这一点，这里我定义了一个数据驱动测试的示例：

\begin{codebox}[python]
# data_driven_test.py
from unittest import TestCase, main
from utils import to_str
class DataDrivenTestCase(TestCase):
    def test_good(self):
        good_cases = [
            (b"my bytes", "my bytes"),
            ("no error", b"no error"),  # This one will fail
            ("other str", "other str"),
            ...
        ]
        for value, expected in good_cases:
            with self.subTest(value):
                self.assertEqual(expected, to_str(value))
    def test_bad(self):
        bad_cases = [
            (object(), TypeError),
            (b"\xfa\xfa", UnicodeDecodeError),
            ...
        ]
        for value, exception in bad_cases:
            with self.subTest(value):
                with self.assertRaises(exception):
                    to_str(value)
if __name__ == "__main__":
    main()
\end{codebox}

\texttt{"no error"} 这个测试用例会失败并打印出有用的错误消息，但所有其他用例仍然会被测试并确认通过：

\begin{codebox}[python]
$ python3 data_driven_test.py
.
===============================================================
FAIL: test_good (__main__.DataDrivenTestCase) [no error]
---------------------------------------------------------------
Traceback (most recent call last):
  File "testing/data_driven_test.py", line 18, in test_good
    self.assertEqual(expected, to_str(value))
AssertionError: b'no error' != 'no error'
---------------------------------------------------------------
Ran 2 tests in 0.001s
FAILED (failures=1)
\end{codebox}

随着项目复杂度和测试需求的增长，你迟早可能会觉得 \texttt{unittest} 及其能力不够用了。一旦到了那一天，开源包 \texttt{pytest}（https://pytest.org）及其大量的社区插件，作为替代的测试运行器（test runner）会格外有用。

\subsection{需要记住的要点}

\begin{itemize}
\item 你可以通过继承内置模块 \texttt{unittest} 中的 \texttt{TestCase} 类、并为每个想要测试的行为定义一个方法来创建测试。\texttt{TestCase} 类上的测试方法必须以 \texttt{test} 开头。
\item 使用 \texttt{TestCase} 类定义的各种辅助方法（例如 \texttt{assertEqual}）来确认测试中的预期行为，而不是使用内置的 \texttt{assert} 语句。
\item 考虑使用 \texttt{subTest} 辅助方法来编写数据驱动测试，以减少样板代码。
\end{itemize}

\section{Item 109: 优先编写集成测试而非单元测试}

软件测试的方法有许多，其范围远超出 Python 本身，包括测试驱动开发、基于属性的测试、变异测试，以及代码与分支覆盖率报告。在 Python 的内置包和社区包中，你能找到编写各种类型和风格的自动化测试的出色工具（参见 Item 116「知道去哪里寻找社区构建的模块」）。因此，在 Python 中的问题不是「能否且应该编写测试」，而是：测试做到什么程度才算足够？测试究竟应当验证什么？

最好把 Python 中的测试看作代码的保险单。好的测试让你确信代码是正确的。如果你重构或扩展代码，那么验证行为（而非实现）的测试能让你轻松看出哪些地方发生了变化。这听起来有悖直觉，但拥有精心编写的测试实际上会让修改 Python 代码更容易，而不是更难。

与其他语言一样，测试可以覆盖代码库的许多不同层次。单元测试（unit test）验证庞大得多的系统中的局部片段。当你有大量边界情况（edge case）并且需要确保一切都被正确处理时，它们很有用。它们运行迅速，因为只用到程序的一小部分。它们通常使用模拟对象（mock）来构建（参见 Item 111「使用 Mock 测试具有复杂依赖的代码」）。

集成测试（integration test）验证多个组件协同工作。它们通常运行更慢、编写更难（示例参见 Item 110「用 setUp、tearDown、setUpModule 与 tearDownModule 隔离各个测试」）。然而，集成测试在 Python 中尤为重要，因为你无法保证子系统真的能够互操作，除非你加以证明（参见 Item 3「切勿指望 Python 在编译期检测错误」）。静态类型（statically typed）语言可以利用类型信息来大致判断组件的契合度，但在动态语言（dynamic language）中，同样借助类型往往要困难得多（参见 Item 124「考虑通过 typing 进行静态分析以消除 Bug」），甚至实际上不可行。

一般而言，在 Python 中最好编写集成测试。但如果你注意到代码的某些部分也有大量边界条件（boundary condition）需要探索，那么为这些行为编写单元测试也值得考虑。你应当避免的，是只写单元测试。例如，假设我在构建用于控制烤面包机的嵌入式软件。这里我定义一个烤面包机类，可以设定「熟度（\texttt{doneness}）」等级、压下面包或弹出烤好的面包：

\begin{codebox}[python]
class Toaster:
    def __init__(self, timer):
        self.timer = timer
        self.doneness = 3
        self.hot = False
    def _get_duration(self):
        return max(0.1, min(120, self.doneness * 10))
    def push_down(self):
        if self.hot:
            return
        self.hot = True
        self.timer.countdown(self._get_duration(), self.pop_up)
    def pop_up(self):
        print("Pop!")  # 松开弹簧
        self.hot = False
        self.timer.end()
\end{codebox}

\texttt{Toaster} 类依赖一个计时器（timer），在面包烤好后将其弹出。这个计时器应当能够被任意多次重置。这里我使用内置的 \texttt{threading} 模块中的 \texttt{Timer} 类来实现：

\begin{codebox}[python]
import threading
class ReusableTimer:
    def __init__(self):
        self.timer = None
    def countdown(self, duration, callback):
        self.end()
        self.timer = threading.Timer(duration, callback)
        self.timer.start()
    def end(self):
        if self.timer:
            self.timer.cancel()
\end{codebox}

定义了这两个类之后，我可以轻松演练烤面包机的功能，证明它能在面包烤焦之前为其加热并弹出：

\begin{codebox}[python]
toaster = Toaster(ReusableTimer())
print("Initially hot:  ", toaster.hot)
toaster.doneness = 5
toaster.push_down()
print("After push down:", toaster.hot)
# 时间流逝
...
print("After time:     ", toaster.hot)
>>>
Initially hot:   False
After push down: True
Pop!
After time:      False
\end{codebox}

如果我想为 \texttt{Toaster} 类编写一个单元测试（unit test），可以用内置的 \texttt{unittest} 模块这样做（参见 Item 108「在 TestCase 子类中验证相关行为」），这里我把 \texttt{ReusableTimer} 类完全用模拟对象（mock）替换掉：

\begin{codebox}[python]
from unittest import TestCase
from unittest.mock import Mock
class ToasterUnitTest(TestCase):
    def test_start(self):
        timer = Mock(spec=ReusableTimer)
        toaster = Toaster(timer)
        toaster.push_down()
        self.assertTrue(toaster.hot)
         timer.countdown.assert_called_once_with(
     30, toaster.pop_up)
    def test_end(self):
        timer = Mock(spec=ReusableTimer)
        toaster = Toaster(timer)
        toaster.hot = True
        toaster.pop_up()
        self.assertFalse(toaster.hot)
        timer.end.assert_called_once()
...
>>>
Pop!
..
---------------------------------------------------------------
Ran 2 tests in 0.000s
OK
\end{codebox}

为 \texttt{ReusableTimer} 类编写单元测试同样可以模拟（mock）其依赖项：

\begin{codebox}[python]
from unittest import mock
class ReusableTimerUnitTest(TestCase):
    def test_countdown(self):
        my_func = lambda: None
        with mock.patch("threading.Timer"):
            timer = ReusableTimer()
            timer.countdown(0.1, my_func)
            threading.Timer.assert_called_once_with(0.1, my_func)
            timer.timer.start.assert_called_once()
    def test_end(self):
        my_func = lambda: None
        with mock.patch("threading.Timer"):
            timer = ReusableTimer()
            timer.countdown(0.1, my_func)
            timer.end()
            timer.timer.cancel.assert_called_once()
...
>>>
..
---------------------------------------------------------------
Ran 2 tests in 0.001s
OK
\end{codebox}

这些单元测试可以工作，但它们需要相当多的搭建工作，还要摆弄模拟对象。相比之下，考虑这个单一的集成测试，它在不使用任何模拟对象的情况下，一起验证 \texttt{Toaster} 与 \texttt{ReusableTimer} 类：

\begin{codebox}[python]
class ToasterIntegrationTest(TestCase):
    def setUp(self):
        self.timer = ReusableTimer()
        self.toaster = Toaster(self.timer)
        self.toaster.doneness = 0
    def test_wait_finish(self):
        self.assertFalse(self.toaster.hot)
        self.toaster.push_down()
        self.assertTrue(self.toaster.hot)
        self.timer.timer.join()
        self.assertFalse(self.toaster.hot)
    def test_cancel_early(self):
        self.assertFalse(self.toaster.hot)
        self.toaster.push_down()
        self.assertTrue(self.toaster.hot)
        self.toaster.pop_up()
        self.assertFalse(self.toaster.hot)
...
>>>
Pop!
.Pop!
.
---------------------------------------------------------------
Ran 2 tests in 0.108s
OK
\end{codebox}

这个测试清晰、简洁，并且聚焦于端到端（end-to-end）行为而非实现细节。也许我对它唯一的抱怨是：为了正确等待 \texttt{threading.Timer} 实例结束，它访问了 \texttt{ReusableTimer} 类的内部（使用 \texttt{join} 方法）。但这就是 Python，而测试时拥有这种访问能力正是这门语言的主要优势之一。

与此单一集成测试相比，前面分别针对 \texttt{Toaster} 和 \texttt{ReusableTimer} 类的单元测试显得冗余且不必要的复杂。不过，单元测试可能为这段代码带来一个潜在好处：测试 \texttt{doneness} 设定的边界，确保它既不会太长也不会太短：

\begin{codebox}[python]
class DonenessUnitTest(TestCase):
    def setUp(self):
        self.toaster = Toaster(ReusableTimer())
    def test_min(self):
        self.toaster.doneness = 0
        self.assertEqual(0.1, self.toaster._get_duration())
    def test_max(self):
        self.toaster.doneness = 1000
        self.assertEqual(120, self.toaster._get_duration())
...
>>>
..
---------------------------------------------------------------
Ran 2 tests in 0.000s
OK
\end{codebox}

这就是你在 Python 中应当编写的测试的正确平衡：一定要有为端到端行为编写的集成测试，或许再为复杂的边界情况（edge case）编写单元测试。大多数时候很容易避免使用模拟对象，只在确有充分理由时才使用它们（参见 Item 112「封装依赖以便于模拟与测试」）。除此之外，别忘了你仍然需要更大的系统测试（system test），来验证你的 Python 程序如何与相应的 Web 客户端、API 端点、移动应用、数据库等交互。

\subsection{需要记住的要点}

\begin{itemize}
\item 集成测试验证多个组件协同工作的行为，而单元测试只验证单个组件自身。
\item 由于 Python 高度动态的特性，集成测试是获得程序正确性信心的绝佳途径——有时甚至是唯一途径。
\item 除集成测试之外，单元测试可用于验证代码库中存在大量边界情况（edge case）或边界条件（boundary condition）的部分。
\end{itemize}

\section{Item 110: 用 setUp、tearDown、setUpModule 与 tearDownModule 隔离各个测试}

\texttt{TestCase} 的子类（见 Item 108「在 \texttt{TestCase} 子类中验证相关行为」）往往需要在测试方法运行前先把测试环境搭建好，这一过程有时称为测试脚手架（test harness）。要做这样的搭建，你可以重写（override）\texttt{TestCase} 父类的 \texttt{setUp} 和 \texttt{tearDown} 方法。这两个方法会分别在各个测试方法之前与之后被调用，从而让你能够确保每次测试都在隔离（isolation）状态下运行——这是良好测试的一项重要最佳实践。

例如，下面我定义了一个 \texttt{TestCase} 子类，它会在每个测试之前创建一个临时目录，并在每个测试结束后删除其中的内容：

\begin{codebox}[python]
# environment_test.py
from pathlib import Path
from tempfile import TemporaryDirectory
from unittest import TestCase, main
class EnvironmentTest(TestCase):
    def setUp(self):
        self.test_dir = TemporaryDirectory()
        self.test_path = Path(self.test_dir.name)
    def tearDown(self):
        self.test_dir.cleanup()
    def test_modify_file(self):
        with open(self.test_path / "data.bin", "w") as f:
            ...
if __name__ == "__main__":
    main()
\end{codebox}

随着程序变得复杂，你可以增加额外的测试来验证各模块之间的端到端（end-to-end）交互，而不只是孤立地测试代码（见 Item 109「优先选择集成测试而非单元测试」）。常见的一个问题是：为集成测试搭建测试环境可能计算开销很大，并且需要耗费大量挂钟时间（wall-clock time）。例如，在运行集成测试之前，你可能需要启动一个数据库进程，并等待它完成索引加载。这种延迟使得在 \texttt{TestCase} 子类的 \texttt{setUp} 和 \texttt{tearDown} 方法中为每个测试都做准备工作与清理变得不切实际。

为了应对这种情况，\texttt{unittest} 模块还支持模块级的测试脚手架初始化。你可以一次性配置好开销昂贵的资源，然后让所有 \texttt{TestCase} 类及其测试方法在不必重复该初始化的前提下运行。之后，当模块中所有测试都执行完毕，测试脚手架也可以被一次性拆除。下面我在包含 \texttt{TestCase} 子类的模块中定义 \texttt{setUpModule} 和 \texttt{tearDownModule} 函数来利用这一行为：

\begin{codebox}[python]
# integration_test.py
from unittest import TestCase, main
def setUpModule():
    print("* Module setup")
def tearDownModule():
    print("* Module clean-up")
class IntegrationTest(TestCase):
    def setUp(self):
        print("* Test setup")
    def tearDown(self):
        print("* Test clean-up")
    def test_end_to_end1(self):
        print("* Test 1")
    def test_end_to_end2(self):
        print("* Test 2")
if __name__ == "__main__":
    main()
$ python3 integration_test.py
* Module setup
* Test setup
* Test 1
* Test clean-up
.* Test setup
* Test 2
* Test clean-up
.* Module clean-up
---------------------------------------------------------------
Ran 2 tests in 0.000s
OK
\end{codebox}

\texttt{setUpModule} 函数由 \texttt{unittest} 只运行一次，且发生在任何 \texttt{setUp} 方法被调用之前。类似地，\texttt{tearDownModule} 发生在 \texttt{tearDown} 方法被调用之后。

\subsection{需要记住的要点}

\begin{itemize}
\item 使用 \texttt{TestCase} 的 \texttt{setUp} 和 \texttt{tearDown} 方法，确保你的各个测试彼此隔离，并保证一个干净的测试环境。
\item 对于集成测试，使用 \texttt{setUpModule} 和 \texttt{tearDownModule} 这两个模块级函数来管理工作所需的测试脚手架，其生命周期覆盖整个测试模块及其包含的所有 \texttt{TestCase} 子类。
\end{itemize}

\section{Item 111: 使用 Mock 测试具有复杂依赖的代码}

编写测试时另一个常见需求（参见 Item 108「在 \texttt{TestCase} 子类中验证相关行为」）是，当真实对象难以使用或运行太慢时，使用 mock 函数和类来模拟行为。例如，假设我需要一个程序来维护动物园里动物的喂食时间表。这里我定义一个函数，查询数据库中某一物种的所有动物，并返回它们最近一次进食的时间：

\begin{codebox}[python]
class DatabaseConnection:
    ...
def get_animals(database, species):
    # 查询数据库
    ...
    # 返回一个由 (name, last_mealtime) 元组组成的列表
\end{codebox}

我该如何获取一个 \texttt{DatabaseConnection} 实例，用来测试这个函数？这里我尝试创建一个实例，并将其传入被测试的函数：

\begin{codebox}[python]
database = DatabaseConnection("localhost", "4444")
get_animals(database, "Meerkat")
>>>
Traceback ...
DatabaseConnectionError: Not connected
\end{codebox}

没有数据库在运行，因此这当然会失败。一种解决办法是在测试中真正启动一个数据库服务器并连接它。然而，为了仅仅运行一个简单的单元测试，要把启动数据库、配置其结构（schema）、填充数据等步骤完全自动化，需要做大量工作。此外，启动数据库服务器可能会消耗大量的实际时间（wall-clock time），这会拖慢这些单元测试，并让它们更难维护（关于一种可能的解决方案，参见 Item 110「用 \texttt{setUp}、\texttt{tearDown}、\texttt{setUpModule} 和 \texttt{tearDownModule} 将测试彼此隔离」）。

另一种办法是把数据库 mock 掉。mock 让你能够在给定一组预期调用的前提下，为依赖函数提供预期的返回值。需要注意的是，不要把 mock 与伪对象（fake）混淆。伪对象（fake）会提供 \texttt{DatabaseConnection} 类的大部分行为，但实现更简单，例如一个基本的、单线程、无持久化的内存数据库。

Python 提供了 \texttt{unittest.mock} 内置模块，用于创建 mock 并在测试中使用它们。这里我定义一个 \texttt{Mock} 实例，用来模拟 \texttt{get\_animals} 函数，而无需真正连接数据库：

\begin{codebox}[python]
from datetime import datetime
from unittest.mock import Mock
mock = Mock(spec=get_animals)
expected = [
    ("Spot", datetime(2024, 6, 5, 11, 15)),
    ("Fluffy", datetime(2024, 6, 5, 12, 30)),
    ("Jojo", datetime(2024, 6, 5, 12, 45)),
]
mock.return_value = expected
\end{codebox}

\texttt{Mock} 类会创建一个 mock 函数。该 mock 的 \texttt{return\_value} 属性就是调用它时要返回的值。\texttt{spec} 参数表明这个 mock 应当表现得像所给的对象（在本例中是一个函数），如果以错误的方式使用它，就会报错。例如，这里我尝试把这个 mock 函数当作一个带有属性的 mock 对象来对待：

\begin{codebox}[python]
mock.does_not_exist
>>>
Traceback ...
AttributeError: Mock object has no attribute 'does_not_exist'
\end{codebox}

一旦创建好，我就可以调用这个 mock、获取其返回值，并验证它返回的内容符合预期。我用了一个独一无二的 \texttt{object()} 值作为 \texttt{database} 参数，因为它实际上不会被 mock 用来做任何事；我关心的只是 \texttt{database} 参数被正确地传递给了任何需要 \texttt{DatabaseConnection} 实例才能工作的依赖函数：

\begin{codebox}[python]
database = object()
result = mock(database, "Meerkat")
assert result == expected
\end{codebox}

这验证了 mock 作出了正确的响应，但我怎么知道调用 mock 的代码是否提供了正确的参数？为此，\texttt{Mock} 类提供了 \texttt{assert\_called\_once\_with} 方法，它能验证确实进行了一次完全匹配给定参数的调用：

\begin{codebox}[python]
mock.assert_called_once_with(database, "Meerkat")
\end{codebox}

如果我提供了错误的参数，就会抛出异常，而任何使用了该断言的 \texttt{TestCase} 都会失败：

\begin{codebox}[python]
mock.assert_called_once_with(database, "Giraffe")
>>>
Traceback ...
AssertionError: expected call not found.
Expected: mock(<object object at 0x104728900>, 'Giraffe')
  Actual: mock(<object object at 0x104728900>, 'Meerkat')
\end{codebox}

如果我其实并不关心某些个别参数（例如具体使用了哪个 \texttt{database} 对象），那么我可以通过 \texttt{unittest.mock.ANY} 常量来表明，某个参数取任何值都是可以的。我还可以使用 \texttt{Mock} 的 \texttt{assert\_called\_with} 方法，来验证对 mock 的最近一次调用（这种情况下也许发生过多次调用）符合我的预期：

\begin{codebox}[python]
from unittest.mock import ANY
mock = Mock(spec=get_animals)
mock("database 1", "Rabbit")
mock("database 2", "Bison")
mock("database 3", "Meerkat")
mock.assert_called_with(ANY, "Meerkat")
\end{codebox}

当一个参数并非被测行为的核心时，\texttt{ANY} 在测试中非常有用。与其过度指定（overspecify）测试、不得不把各种测试参数的预期都传递下去，不如更宽松地使用 \texttt{ANY}、倾向于把测试指定得少一些，这往往是更值得的做法。

\texttt{Mock} 类还能轻松地模拟抛出异常：

\begin{codebox}[python]
class MyError(Exception):
    pass
mock = Mock(spec=get_animals)
mock.side_effect = MyError("Whoops! Big problem")
result = mock(database, "Meerkat")
>>>
Traceback ...
MyError: Whoops! Big problem
\end{codebox}

还有更多特性可用，因此请务必查看该模块的文档，了解全部选项（\texttt{https://docs.python.org/3/library/unittest.mock.html}）。

既然我已经展示了 \texttt{Mock} 对象的工作原理，我就可以把它应用到一个真实的测试场景中，展示如何在编写测试时有效地使用它。这里我定义一组与数据库交互的函数，用来完成动物园里动物的喂食巡查：

\begin{codebox}[python]
def get_food_period(database, species):
    # 查询数据库
    ...
    # 返回一个时间差
def feed_animal(database, name, when):
    # 写入数据库
    ...
def do_rounds(database, species):
    now = datetime.now()
    feeding_timedelta = get_food_period(database, species)
    animals = get_animals(database, species)
    fed = 0
    for name, last_mealtime in animals:
        if (now - last_mealtime) > feeding_timedelta:
            feed_animal(database, name, now)
            fed += 1
    return fed
\end{codebox}

我的测试目标是：验证当 \texttt{do\_rounds} 运行时，正确的动物得到了喂食，最新的喂食时间被记录到了数据库，并且函数返回的实际喂食动物总数与正确的总数相符。为了做到这些，我需要把 \texttt{datetime.now} mock 掉，这样我的测试就能期待一个稳定的时间，不受程序执行时刻的影响。我需要把 \texttt{get\_food\_period} 和 \texttt{get\_animals} mock 掉，让它们返回本应来自数据库的值。我还需要把 \texttt{feed\_animal} mock 掉，让它接收本应写回数据库的数据。

问题是：即便我知道如何创建这些 mock 函数并设置预期，我又如何让被测试的 \texttt{do\_rounds} 函数使用 mock 依赖函数，而不是真实版本？一种办法是把所有东西都作为仅限关键字（keyword-only）参数注入（背景知识参见 Item 37「用仅限关键字和仅限位置参数增强清晰度」）：

\begin{codebox}[python]
def do_rounds(
    database,
    species,
    *,
    now_func=datetime.now,
    food_func=get_food_period,
    animals_func=get_animals,
    feed_func=feed_animal
):
    now = now_func()
    feeding_timedelta = food_func(database, species)
    animals = animals_func(database, species)
    fed = 0
    for name, last_mealtime in animals:
        if (now - last_mealtime) > feeding_timedelta:
            feed_func(database, name, now)
            fed += 1
    return fed
\end{codebox}

为了测试这个函数，我需要预先创建所有的 \texttt{Mock} 实例，并设置它们的预期：

\begin{codebox}[python]
from datetime import timedelta
now_func = Mock(spec=datetime.now)
now_func.return_value = datetime(2024, 6, 5, 15, 45)
food_func = Mock(spec=get_food_period)
food_func.return_value = timedelta(hours=3)
animals_func = Mock(spec=get_animals)
animals_func.return_value = [
    ("Spot", datetime(2024, 6, 5, 11, 15)),
    ("Fluffy", datetime(2024, 6, 5, 12, 30)),
    ("Jojo", datetime(2024, 6, 5, 12, 45)),
]
feed_func = Mock(spec=feed_animal)
\end{codebox}

然后，我可以通过把这些 mock 传入 \texttt{do\_rounds} 函数，来覆盖默认参数，从而运行测试：

\begin{codebox}[python]
result = do_rounds(
    database,
    "Meerkat",
    now_func=now_func,
    food_func=food_func,
    animals_func=animals_func,
    feed_func=feed_func,
)
assert result == 2
\end{codebox}

最后，我可以验证对依赖函数的所有调用都符合我的预期：

\begin{codebox}[python]
from unittest.mock import call
food_func.assert_called_once_with(database, "Meerkat")
animals_func.assert_called_once_with(database, "Meerkat")
feed_func.assert_has_calls(
    [
        call(database, "Spot", now_func.return_value),
        call(database, "Fluffy", now_func.return_value),
    ],
    any_order=True,
)
\end{codebox}

我没有验证 \texttt{datetime.now} 这个 mock 的参数，也没有验证它的调用次数，因为这一点已被函数的返回值间接验证。对于 \texttt{get\_food\_period} 和 \texttt{get\_animals}，我使用 \texttt{assert\_called\_once\_with} 来验证它们各被调用了一次，且参数符合指定值。对于 \texttt{feed\_animal} 函数，我使用 \texttt{unittest.mock.call} 辅助函数和 \texttt{assert\_has\_calls} 方法，验证确实进行了两次写入数据库的调用——而且调用的顺序无关紧要。

这种使用仅限关键字参数来注入 mock 的做法是可行的，但它相当繁琐，而且需要修改每一个你想测试的函数。\texttt{unittest.mock.patch} 系列函数让注入 mock 变得更容易。它会临时重新指派一个模块或类的属性，比如我上面定义的那些访问数据库的函数。例如，这里我使用 \texttt{patch} 把 \texttt{get\_animals} 替换成一个 mock：

\begin{codebox}[python]
from unittest.mock import patch
print("Outside patch:", get_animals)
with patch("__main__.get_animals"):
    print("Inside patch: ", get_animals)
print("Outside again:", get_animals)
>>>
Outside patch: <function get_animals at 0x104eda160>
Inside patch:  <MagicMock name='get_animals' id='4397863264'>
Outside again: <function get_animals at 0x104eda160>
\end{codebox}

\texttt{patch} 适用于许多模块、类和属性。它可以用在 \texttt{with} 语句中（参见 Item 82「考虑用 \texttt{contextlib} 和 \texttt{with} 语句实现可复用的 \texttt{try}/\texttt{finally} 行为」），用作函数装饰器（参见 Item 38「用 \texttt{functools.wraps} 定义函数装饰器」），或者用在 \texttt{TestCase} 类的 \texttt{setUp} 和 \texttt{tearDown} 方法中（参见 Item 110「用 \texttt{setUp}、\texttt{tearDown}、\texttt{setUpModule} 和 \texttt{tearDownModule} 将测试彼此隔离」）。

不过，\texttt{patch} 并非在所有情况下都有效。例如，为了测试 \texttt{do\_rounds}，我需要把 \texttt{datetime.now} 类方法返回的当前时间 mock 掉。Python 不允许我这样做，因为 \texttt{datetime} 类定义在一个 C 扩展模块中，无法以这种方式被修改：

\begin{codebox}[python]
fake_now = datetime(2024, 6, 5, 15, 45)
with patch("datetime.datetime.now"):
    datetime.now.return_value = fake_now
>>>
Traceback ...
TypeError: cannot set 'now' attribute of immutable type
'datetime.datetime'
The above exception was the direct cause of the following
exception:
Traceback ...
TypeError: cannot set 'now' attribute of immutable type
'datetime.datetime'
\end{codebox}

为了绕过这一点，我可以再创建一个用于获取时间的辅助函数，它可以被 \texttt{patch}：

\begin{codebox}[python]
def get_do_rounds_time():
    return datetime.now()
def do_rounds(database, species):
    now = get_do_rounds_time()
    ...
with patch("__main__.get_do_rounds_time"):
    ...
\end{codebox}

或者，我可以为 \texttt{datetime.now} 这个 mock 使用一个仅限关键字参数，而对其他所有 mock 使用 \texttt{patch}：

\begin{codebox}[python]
def do_rounds(database, species, *, now_func=datetime.now):
    now = now_func()
    feeding_timedelta = get_food_period(database, species)
    animals = get_animals(database, species)
    fed = 0
    for name, last_mealtime in animals:
        if (now - last_mealtime) > feeding_timedelta:
            feed_animal(database, name, now)
            fed += 1
    return fed
\end{codebox}

我打算采用后一种办法。现在，我可以使用 \texttt{patch.multiple} 函数来创建多个 mock，然后设置它们的预期：

\begin{codebox}[python]
from unittest.mock import DEFAULT
with patch.multiple(
    "__main__",
    autospec=True,
    get_food_period=DEFAULT,
    get_animals=DEFAULT,
    feed_animal=DEFAULT,
):
    now_func = Mock(spec=datetime.now)
    now_func.return_value = datetime(2024, 6, 5, 15, 45)
    get_food_period.return_value = timedelta(hours=3)
    get_animals.return_value = [
        ("Spot", datetime(2024, 6, 5, 11, 15)),
        ("Fluffy", datetime(2024, 6, 5, 12, 30)),
        ("Jojo", datetime(2024, 6, 5, 12, 45)),
    ]
\end{codebox}

\texttt{patch.multiple} 的关键字参数对应于我想在测试期间覆盖的 \texttt{\_\_main\_\_} 模块中的名称。\texttt{DEFAULT} 值表示我想为每个名称创建一个标准的 \texttt{Mock} 实例。由于 \texttt{autospec=True} 参数，所有生成的 mock 都会遵循它们所模拟对象的规范。

准备工作就绪后，我就可以运行测试，并在使用了 \texttt{patch.multiple} 的 \texttt{with} 语句内部，验证相关调用是否正确：

\begin{codebox}[python]
    result = do_rounds(database, "Meerkat", now_func=now_func)
    assert result == 2
    get_food_period.assert_called_once_with(database, "Meerkat")
    get_animals.assert_called_once_with(database, "Meerkat")
    feed_animal.assert_has_calls(
        [
            call(database, "Spot", now_func.return_value),
            call(database, "Fluffy", now_func.return_value),
        ],
        any_order=True,
    )
\end{codebox}

这些 mock 按预期工作，但重要的是要认识到，通过重构代码使其在设计中更易测试，可以进一步提高这些测试的可读性并减少样板代码（boilerplate）（参见 Item 112「封装依赖以便于 Mock 和测试」）。

\subsection{需要记住的要点}

\begin{itemize}
\item \texttt{unittest.mock} 模块提供了一种借助 \texttt{Mock} 类来模拟接口行为的方法。当难以搭建被测试代码所需的依赖时，mock 在测试中非常有用。
\item 使用 mock 时，既要验证被测试代码的行为，也要验证依赖函数被该代码调用的方式，这一点很重要，可以借助 \texttt{Mock.assert\_called\_once\_with} 系列方法来实现。
\item 仅限关键字参数与 \texttt{unittest.mock.patch} 系列函数，都可以用来向被测试代码中注入 mock。
\end{itemize}

\section{Item 112: 封装依赖以便于模拟与测试}

在上一节（见 Item 111「用模拟对象测试带有复杂依赖的代码」）中，我展示了如何使用内置模块 \texttt{unittest.mock}（包括 \texttt{Mock} 类与 \texttt{patch} 系列函数）来编写带有复杂依赖（例如数据库）的测试。然而，由此产生的测试代码需要大量样板代码（boilerplate），这会让代码的新读者更难理解这些测试到底想要验证什么。

改进这些测试的一种方法是使用一个包装对象（wrapper object）来封装数据库的接口，而不是把一个 \texttt{DatabaseConnection} 对象作为参数传递给函数。重构代码以使用更好的抽象（见 Item 123「考虑用 \texttt{warnings} 重构并迁移用法」了解其中一种做法）通常是值得的，因为它有助于创建模拟对象并编写测试。这里我把上一节中的各种数据库辅助函数重新定义为类的方法，而不是独立的函数：

\begin{codebox}[python]
class ZooDatabase:
    ...
    def get_animals(self, species):
        ...
    def get_food_period(self, species):
        ...
    def feed_animal(self, name, when):
        ...
\end{codebox}

现在我可以重新定义 \texttt{do\_rounds} 函数，让它调用 \texttt{ZooDatabase} 对象上的方法：

\begin{codebox}[python]
from datetime import datetime
def do_rounds(database, species, *, now_func=datetime.now):
    now = now_func()
    feeding_timedelta = database.get_food_period(species)
    animals = database.get_animals(species)
    fed = 0
    for name, last_mealtime in animals:
        if (now - last_mealtime) >= feeding_timedelta:
            database.feed_animal(name, now)
            fed += 1
    return fed
\end{codebox}

现在为 \texttt{do\_rounds} 编写测试要容易得多，因为我不再需要用 \texttt{unittest.mock.patch} 把模拟对象注入到被测代码中。相反，我可以创建一个 \texttt{Mock} 实例来代表一个 \texttt{ZooDatabase}，并把它作为 \texttt{database} 参数传入。\texttt{Mock} 类会为任何被访问的属性名返回一个模拟对象。这些属性可以像方法一样被调用，接着我就可以用它们来设定预期并验证调用。这让模拟一个类的所有方法变得很容易：

\begin{codebox}[python]
from unittest.mock import Mock
database = Mock(spec=ZooDatabase)
print(database.feed_animal)
database.feed_animal()
database.feed_animal.assert_any_call()
>>>
<Mock name='mock.feed_animal' id='4386901024'>
\end{codebox}

我可以用 \texttt{ZooDatabase} 封装的方式重写 \texttt{Mock} 的建立代码：

\begin{codebox}[python]
from datetime import timedelta
from unittest.mock import call
now_func = Mock(spec=datetime.now)
now_func.return_value = datetime(2019, 6, 5, 15, 45)
database = Mock(spec=ZooDatabase)
database.get_food_period.return_value = timedelta(hours=3)
database.get_animals.return_value = [
    ("Spot", datetime(2019, 6, 5, 11, 15)),
    ("Fluffy", datetime(2019, 6, 5, 12, 30)),
    ("Jojo", datetime(2019, 6, 5, 12, 55)),
]
\end{codebox}

然后我可以运行被测函数，并验证所有依赖方法都按预期被调用：

\begin{codebox}[python]
result = do_rounds(database, "Meerkat", now_func=now_func)
assert result == 2
database.get_food_period.assert_called_once_with("Meerkat")
database.get_animals.assert_called_once_with("Meerkat")
database.feed_animal.assert_has_calls(
    [
        call("Spot", now_func.return_value),
        call("Fluffy", now_func.return_value),
    ],
    any_order=True,
)
\end{codebox}

在使用 \texttt{Mock} 模拟类时，使用 \texttt{spec} 参数特别有用，因为它能确保被测代码不会意外调用一个拼写错误的方法名。这让你能避免一个常见的陷阱：同样的 bug 同时存在于代码和单元测试中，从而掩盖一个真正的错误，而这个错误之后会在生产环境中暴露出来：

\begin{codebox}[python]
database.bad_method_name()
>>>
Traceback ...
AttributeError: Mock object has no attribute 'bad_method_name'
\end{codebox}

如果我想用一个中等层级的集成测试（integration test，见 Item 109「优先使用集成测试而非单元测试」）来端到端地测试这个程序，我仍然需要一种方法把模拟的 \texttt{ZooDatabase} 注入到程序中。我可以创建一个充当依赖注入（dependency injection）接缝（seam）的辅助函数来实现这一点。这里我定义这样一个辅助函数，它通过 \texttt{global} 语句（见 Item 120「考虑用模块作用域代码配置部署环境」了解背景）在模块作用域（module scope）中缓存一个 \texttt{ZooDatabase} 对象：

\begin{codebox}[python]
DATABASE = None
def get_database():
    global DATABASE
    if DATABASE is None:
        DATABASE = ZooDatabase()
    return DATABASE
def main(argv):
    database = get_database()
    species = argv[1]
    count = do_rounds(database, species)
    print(f"Fed {count} {species}(s)")
    return 0
\end{codebox}

现在我可以用 \texttt{patch} 注入模拟的 \texttt{ZooDatabase}，运行测试，并验证程序的输出。这里我没有使用模拟的 \texttt{datetime.now}；相反，我依赖模拟对象返回的数据库记录相对于当前时间，以产生与单元测试类似的行为。这种方法比把一切都模拟掉更不稳定（flaky），但它也测试了更大的覆盖面（surface area）：

\begin{codebox}[python]
import contextlib
import io
from unittest.mock import patch
with patch("__main__.DATABASE", spec=ZooDatabase):
    now = datetime.now()
    DATABASE.get_food_period.return_value = timedelta(hours=3)
    DATABASE.get_animals.return_value = [
        ("Spot", now - timedelta(minutes=4.5)),
        ("Fluffy", now - timedelta(hours=3.25)),
        ("Jojo", now - timedelta(hours=3)),
    ]
    fake_stdout = io.StringIO()
    with contextlib.redirect_stdout(fake_stdout):
        main(["program name", "Meerkat"])
    found = fake_stdout.getvalue()
    expected = "Fed 2 Meerkat(s)\n"
    assert found == expected
\end{codebox}

结果与我的预期一致。创建这个集成测试是直截了当的，因为我从设计实现时就让它更易于测试。

\subsection{需要记住的要点}

\begin{itemize}
\item 当单元测试需要大量重复的样板代码来建立模拟对象时，一种解决办法是把依赖的功能封装到更容易被模拟的类中。
\item 内置模块 \texttt{unittest.mock} 的 \texttt{Mock} 类通过为每次被访问的属性返回一个新模拟对象（它可以充当模拟方法）来模拟类。
\item 对于端到端测试，重构代码以拥有更多辅助函数是很有价值的，这些辅助函数可以充当显式的接缝（seam），用于在测试中注入模拟依赖。
\end{itemize}

\section{Item 113: 用 assertAlmostEqual 控制浮点数测试的精度}

Python 的 \texttt{float} 类型是双精度浮点数（遵循 IEEE 754 标准）。这种方案有其局限性（见 Item 106「当精度至关重要时使用 \texttt{decimal}」），但浮点数在许多场景下都很有用，并且在 Python 中得到了良好支持。

通常，测试数学代码以检查边界条件（boundary conditions）和其他潜在的错误来源是很重要的（详见 Item 109「优先使用集成测试而非单元测试」）。遗憾的是，编写涉及浮点数的自动化测试可能很棘手。例如，这里我使用内置模块 \texttt{unittest} 定义一个测试，它试图（并且失败地）验证表达式 \texttt{5 / 3} 的结果：

\begin{codebox}[python]
import unittest
class MyTestCase(unittest.TestCase):
    def test_equal(self):
        n = 5
        d = 3
        self.assertEqual(1.667, n / d)  # Raises
...
>>>
Traceback ...
AssertionError: 1.667 != 1.6666666666666667
\end{codebox}

问题在于，在 Python 中表达式 \texttt{5 / 3} 得到的是一个无法被精确表示为 \texttt{float} 值的数（小数点之后不断重复的 \texttt{6} 就证明了这一点）。传给 \texttt{assertEqual} 的期望值 \texttt{1.667} 精度不足以精确匹配计算结果。（两者相差 0.000333……）因此，\texttt{assertEqual} 的方法调用失败了。我可以通过让期望值更精确（例如字面量 \texttt{1.6666666666666667}）来解决这个问题。但在实践中，使用这种级别的精度会让数值测试难以维护。由于舍入（rounding）行为，运算顺序的改变就可能产生不同的结果。架构差异（例如 x86 与 AArch64）也有可能影响到结果。

这里我通过重排一个计算来展示这种舍入问题——这种重排看起来不应该影响结果，但实际上却影响了（注意最后一位数字）：

\begin{codebox}[python]
print(5 / 3 * 0.1)
print(0.1 * 5 / 3)
>>>
0.16666666666666669
0.16666666666666666
\end{codebox}

为了在自动化测试中处理这个问题，可以使用 \texttt{TestCase} 类中的 \texttt{assertAlmostEqual} 辅助方法来对浮点数做近似比较。它能正确处理无穷大（infinity）和 NaN 条件，并尽量减小因舍入而引入的误差。这里我用这个方法来验证：当把数字四舍五入到小数点后两位时，它们相等：

\begin{codebox}[python]
class MyTestCase2(unittest.TestCase):
    def test_equal(self):
        ...
        # Changed
        self.assertAlmostEqual(1.667, n / d, places=2)
...
>>>
.
---------------------------------------------------------------
Ran 1 test in 0.000s
OK
\end{codebox}

\texttt{assertAlmostEqual} 的 \texttt{places} 参数在验证小数部分介于 0 和 1 之间的数字时表现良好。但浮点行为和无限循环小数（repeating decimals）也可能影响更大的数字。例如，考虑下面两次计算之间在绝对值上巨大的差异——即便唯一的改动只是给一个系数加上了 \texttt{0.001}：

\begin{codebox}[python]
print(1e24 / 1.1e16)
print(1e24 / 1.101e16)
>>>
90909090.9090909
90826521.34423251
\end{codebox}

这些数值之间的差异约为 82,569。

根据用例的不同，这个误差范围可能重要，也可能不重要。为了让你能够表达对不精确性的容忍度，你可以给 \texttt{assertAlmostEqual} 辅助方法提供一个 \texttt{delta} 参数。这个参数会让方法考量数值之间的绝对差值，并且仅当该差值大于所提供的 \texttt{delta} 时才抛出 \texttt{AssertionError} 异常。

这里我使用这个选项来指定 100,000 的容忍度，它大于 82,569 的差异，从而让两个断言都能通过：

\begin{codebox}[python]
class MyTestCase3(unittest.TestCase):
    def test_equal(self):
        a = 1e24 / 1.1e16
        b = 1e24 / 1.101e16
        self.assertAlmostEqual(90.9e6, a, delta=0.1e6)
        self.assertAlmostEqual(90.9e6, b, delta=0.1e6)
\end{codebox}

在某些情况下，你可能需要断言相反的情况：在给定容忍度或小数位数的前提下，两个数字彼此并不接近。\texttt{TestCase} 类还提供了 \texttt{assertNotAlmostEqual} 方法来让这件事变得简单。为了在测试代码内外比较数字时处理更复杂的用例，内置模块 \texttt{math} 提供了 \texttt{isclose} 函数，它具有类似、甚至更多的功能。

\subsection{需要记住的要点}

\begin{itemize}
\item 由于舍入行为，浮点数（尤其是其小数部分）可能会因所应用的运算顺序不同而改变。
\item 用 \texttt{assertEqual} 测试浮点值可能导致不稳定的测试（flaky tests），因为该方法会考量被比较数字的全部精度。
\item \texttt{assertAlmostEqual} 与 \texttt{assertNotAlmostEqual} 方法允许你指定 \texttt{places} 或 \texttt{delta} 参数，以表明你在比较浮点数时对差异的容忍度。
\end{itemize}

\section{Item 114: 考虑使用 pdb 进行交互式调试}

开发程序时，每个人的代码都会遇到 bug。使用 \texttt{print} 函数有助于追踪许多问题的根源（见 Item 12「理解打印对象时 \texttt{repr} 与 \texttt{str} 的区别」）。针对导致麻烦的特定情况编写测试，是定位问题的另一种好方法（见 Item 109「优先选择集成测试而非单元测试」）。

但这些工具不足以找出所有根本原因。当你需要更强大的工具时，就应当尝试 Python 的内置（built-in）交互式调试器（debugger）。调试器让你可以检查运行中程序的状态、打印局部变量，并逐条语句地单步执行。

在大多数其他编程语言中，使用调试器的方式是先指定源文件里想要停下来的行，然后再执行程序。相比之下，在 Python 中，使用调试器最简单的方式是修改你的程序，在你认为可能出现值得调查的问题之前直接启动调试器。这意味着，为了运行调试器而启动 Python 程序，与正常启动它并没有区别。

要启动调试器，你只需调用 \texttt{breakpoint} 内置（built-in）函数即可。这等价于导入 \texttt{pdb} 内置模块并运行其 \texttt{set\_trace} 函数：

\begin{codebox}[python]
# always_breakpoint.py
import math
def compute_rmse(observed, ideal):
    total_err_2 = 0
    count = 0
    for got, wanted in zip(observed, ideal):
        err_2 = (got - wanted) ** 2
        breakpoint()  # 在此启动调试器
        total_err_2 += err_2
        count += 1
    mean_err = total_err_2 / count
    rmse = math.sqrt(mean_err)
    return rmse
result = compute_rmse(
    [1.8, 1.7, 3.2, 6],
    [2, 1.5, 3, 5],
)
print(result)
\end{codebox}

一旦 \texttt{breakpoint} 函数被执行，程序就会暂停在 \texttt{breakpoint} 调用之后的那行代码之前。启动该程序的终端会变成一个 Python 调试 shell：

\begin{codebox}[python]
$ python3 always_breakpoint.py
> always_breakpoint.py(12)compute_rmse()
-> total_err_2 += err_2
(Pdb)
\end{codebox}

在 \texttt{(Pdb)} 提示符下，你可以输入局部变量的名称来查看它们被打印出的值（或使用 \texttt{p <name>}）。调用 \texttt{locals} 内置函数即可查看所有局部变量的列表。你可以导入模块、检查全局状态、构造新对象，甚至修改运行中程序的某些部分。这个调试提示符并不支持某些 Python 语句和语言特性，但你可以通过 \texttt{interact} 命令进入一个能够访问程序状态的标准 Python REPL（交互式解释器）。

此外，调试器还提供多种特殊命令来控制和理解程序的执行；输入 \texttt{help} 可查看完整列表。

下面三个非常有用的命令让检查运行中的程序变得更容易：

\begin{itemize}
\item \texttt{where}：打印当前执行的调用栈（call stack）。这让你弄清楚自己处于程序的什么位置，以及是如何到达触发断点的地方的。
\item \texttt{up}：将当前作用域沿着执行调用栈上移一级，到达当前函数的调用者。这样你就能检查导致断点的、更高层级程序中的局部变量。
\item \texttt{down}：将当前作用域沿执行调用栈下移一级。
\end{itemize}

当你检查完当前状态后，可以使用下面这五个调试器命令，以不同方式控制程序的执行：

\begin{itemize}
\item \texttt{step}：运行程序直到下一条执行代码行，然后将控制权交还调试器提示符。如果下一条要执行的代码包含函数调用，调试器会停在所调用函数内部。
\item \texttt{next}：运行程序直到当前函数中的下一条执行代码行，然后将控制权交还调试器提示符。如果下一条要执行的代码包含函数调用，调试器会一直运行到被调用函数返回后才会停下。
\item \texttt{return}：运行程序直到当前函数返回，然后将控制权交还调试器提示符。
\item \texttt{continue}：继续运行程序，直到命中下一个断点（无论是通过显式调用 \texttt{breakpoint}，还是遇到由先前调试器命令添加的断点）。
\item \texttt{quit}：退出调试器并结束程序。如果你已找到问题、走得太远，或需要修改程序后重试，就执行这个命令。
\end{itemize}

\texttt{breakpoint} 函数可以在程序中的任何位置调用。如果你知道要调试的问题只在某些特殊情况下才出现，那么你可以编写普通的 Python 代码，在特定条件满足后再调用 \texttt{breakpoint}。例如，这里我仅在某个数据点的平方误差大于 \texttt{1} 时才启动调试器：

\begin{codebox}[python]
# conditional_breakpoint.py
def compute_rmse(observed, ideal):
    ...
    for got, wanted in zip(observed, ideal):
        err_2 = (got - wanted) ** 2
        if err_2 >= 1:  # 若条件为真则启动调试器
            breakpoint()
        total_err_2 += err_2
        count += 1
    ...
result = compute_rmse(
    [1.8, 1.7, 3.2, 7],
    [2, 1.5, 3, 5],
)
print(result)
\end{codebox}

当我运行程序并进入调试器后，可以通过检查局部变量来确认该条件确实为真：

\begin{codebox}[python]
$ python3 conditional_breakpoint.py
> conditional_breakpoint.py(14)compute_rmse()
-> total_err_2 += err_2
(Pdb) wanted
5
(Pdb) got
7
(Pdb) err_2
4
\end{codebox}

另一种进入调试器提示符的实用方式是利用事后（postmortem）调试。它让你能在程序抛出异常并崩溃之后再去调试它。当你不太确定应该把 \texttt{breakpoint} 函数调用放在何处时，这种做法尤其有用。这里我有一个脚本，它会因为某个函数参数中出现了 \texttt{7j} 这个复数而崩溃：

\begin{codebox}[python]
# postmortem_breakpoint.py
import math
def compute_rmse(observed, ideal):
    ...
result = compute_rmse(
    [1.8, 1.7, 3.2, 7j],  # Bad input
    [2, 1.5, 3, 5],
)
print(result)
\end{codebox}

我使用命令行 \texttt{python3 -m pdb -c continue <program path>} 在 \texttt{pdb} 模块的控制下运行程序。\texttt{continue} 命令告诉 \texttt{pdb} 立即启动程序。一旦它开始运行，程序就会遇到问题并自动进入交互式调试器，此时我就可以检查程序状态：

\begin{codebox}[python]
$ python3 -m pdb -c continue postmortem_breakpoint.py
Traceback (most recent call last):
  File "pdb.py", line 1944, in main
    pdb._run(target)
  File "pdb.py", line 1738, in _run
    self.run(target.code)
  File "bdb.py", line 606, in run
    exec(cmd, globals, locals)
  File "<string>", line 1, in <module>
  File "postmortem_breakpoint.py", line 22, in <module>
    result = compute_rmse(
        [1.8, 1.7, 3.2, 7j],  # Bad input
        [2, 1.5, 3, 5],
    )
  File "postmortem_breakpoint.py", line 17, in compute_rmse
    rmse = math.sqrt(mean_err)
TypeError: must be real number, not complex
Uncaught exception. Entering post mortem debugging
Running 'cont' or 'step' will restart the program
> postmortem_breakpoint.py(17)compute_rmse()
-> rmse = math.sqrt(mean_err)
(Pdb) mean_err
(-5.97-17.5j)
\end{codebox}

在交互式 Python 解释器中遇到未捕获异常后，你也可以调用 \texttt{pdb} 模块的 \texttt{pm} 函数来进行事后调试（通常写为一行 \texttt{import pdb; pdb.pm()}）：

\begin{codebox}[python]
$ python3
>>> import my_module
>>> my_module.compute_stddev([5])
Traceback (most recent call last):
  File "<stdin>", line 1, in <module>
  File "my_module.py", line 20, in compute_stddev
    variance = compute_variance(data)
               ^^^^^^^^^^^^^^^^^^^^^^
  File "my_module.py", line 15, in compute_variance
    variance = err_2_sum / (len(data) - 1)
               ~~~~~~~~~~^~~~~~~~~~~~~~~~~
ZeroDivisionError: float division by zero
>>> import pdb; pdb.pm()
> my_module.py(15)compute_variance()
-> variance = err_2_sum / (len(data) - 1)
(Pdb) err_2_sum
0.0
(Pdb) len(data)
1
\end{codebox}

\subsection{需要记住的要点}

\begin{itemize}
\item 你可以在程序中直接调用 \texttt{breakpoint} 内置函数，在你感兴趣的位置启动 Python 交互式调试器。
\item \texttt{pdb} 的 shell 命令让你可以精确控制程序的执行，并让你能在检查程序状态与推进程序执行之间来回切换。
\item \texttt{pdb} 模块既可用于调试独立 Python 程序（使用 \texttt{python -m pdb -c continue <program path>}）中发生的异常，也可用于调试交互式 Python 解释器（使用 \texttt{import pdb; pdb.pm()}）中发生的异常。
\end{itemize}

\section{Item 115: 使用 tracemalloc 了解内存使用与泄漏}

默认实现的 Python（即 CPython）中的内存管理使用引用计数（reference counting）。这保证了只要对一个对象的所有引用都失效，该被引用的对象就会立即从内存中清除，把那块空间腾出来供其他数据使用。CPython 还内置了一个循环检测器（cycle detector），以确保自引用对象最终会被垃圾回收（GC）。

理论上，这意味着大多数 Python 开发者不必担心程序中的内存分配与释放——这些由语言和 CPython 运行时（runtime）自动处理。然而在实践中，当那些已不再有用的引用仍被持有时，程序最终还是会耗尽内存。弄清楚 Python 程序在哪里占用或泄漏内存，可能颇具挑战性。

调试内存使用的一种方式是请 \texttt{gc} 内置（built-in）模块列出垃圾回收器当前已知的所有对象。尽管这是个相当粗糙的工具，但这种方法能让你快速了解程序的内存都用在了何处。这里我定义一个通过持有引用来占满内存的模块：

\begin{codebox}[python]
# waste_memory.py
import os
class MyObject:
    def __init__(self):
        self.data = os.urandom(100)
def get_data():
    values = []
    for _ in range(100):
        obj = MyObject()
        values.append(obj)
    return values
def run():
    deep_values = []
    for _ in range(100):
        deep_values.append(get_data())
    return deep_values
\end{codebox}

然后我运行一个程序，使用 \texttt{gc} 内置模块打印出执行期间创建了多少对象，以及一小部分已分配对象的样本：

\begin{codebox}[python]
# using_gc.py
import gc
found_objects = gc.get_objects()
print("Before:", len(found_objects))
import waste_memory
hold_reference = waste_memory.run()
found_objects = gc.get_objects()
print("After: ", len(found_objects))
for obj in found_objects[:3]:
    print(repr(obj)[:100])
>>>
Before: 6207
After:  16801
<waste_memory.MyObject object at 0x10390aeb8>
<waste_memory.MyObject object at 0x10390aef0>
<waste_memory.MyObject object at 0x10390af28>
...
\end{codebox}

\texttt{gc.get\_objects} 的问题在于，它完全不会告诉你这些对象是如何被分配的。在复杂的程序中，某个特定类的对象可能通过许多不同的方式被分配。比起知道对象的总数，找出究竟是哪段代码负责分配了正在泄漏内存的对象要重要得多。

Python 3.4 引入了一个全新的 \texttt{tracemalloc} 内置（built-in）模块来解决这一问题。\texttt{tracemalloc} 能够将对象回溯到它被分配的位置。使用方式是对内存使用情况进行前后两次快照（snapshot），然后比较它们以查看发生了什么变化。

\begin{codebox}[python]
# top_n.py
import tracemalloc
tracemalloc.start(10)                      # 设置栈深度
time1 = tracemalloc.take_snapshot()        # 之前快照
import waste_memory
x = waste_memory.run()                     # 待调试的使用情况
time2 = tracemalloc.take_snapshot()        # 之后快照
stats = time2.compare_to(time1, "lineno")  # 比较快照
for stat in stats[:3]:
    print(stat)
>>>
waste_memory.py:5: size=2392 KiB (+2392 KiB), count=29994
'(+29994), average=82 B
waste_memory.py:10: size=547 KiB (+547 KiB), count=10001
'(+10001), average=56 B
waste_memory.py:11: size=82.8 KiB (+82.8 KiB), count=100
'(+100), average=848 B
\end{codebox}

输出中的 \texttt{size} 和 \texttt{count} 标签让你立刻就能看清哪些对象主导了程序的内存使用，以及它们在源代码中的何处被分配。

\texttt{tracemalloc} 模块还能打印出每次分配（allocation）的完整调用栈（stack trace）（上限为传给 \texttt{tracemalloc.start} 函数的帧数）。这里我打印出程序中内存占用最大的来源的调用栈：

\begin{codebox}[python]
# with_trace.py
import tracemalloc
tracemalloc.start(10)
time1 = tracemalloc.take_snapshot()
import waste_memory
x = waste_memory.run()
time2 = tracemalloc.take_snapshot()
stats = time2.compare_to(time1, "traceback")
top = stats[0]
print("Biggest offender is:")
print("\n".join(top.traceback.format()))
>>>
Biggest offender is:
  File "with_trace.py", line 11
    x = waste_memory.run()
  File "waste_memory.py", line 20
    deep_values.append(get_data())
  File "waste_memory.py", line 12
    obj = MyObject()
  File "waste_memory.py", line 6
    self.data = os.urandom(100)
\end{codebox}

像这样的调用栈对于搞清楚程序中某个常见函数或类的哪一次具体使用导致了内存消耗，价值最大。

对于更高级的内存分析（profiling）需求，还可以考虑一些社区（community）包（见 Item 116「了解在哪里能找到社区构建的模块」），例如 Memray（https://github.com/bloomberg/memray）。

\subsection{需要记住的要点}

\begin{itemize}
\item 理解 Python 程序如何使用并泄漏内存，可能相当困难。
\item \texttt{gc} 模块能帮你了解存在哪些对象，但它没有任何关于这些对象是如何被分配的信息。
\item \texttt{tracemalloc} 内置模块提供了解内存使用来源的强大工具。
\end{itemize}
